% =============================================================================
% Guion de estudio — presentación a Gonzalo
% Compilar: pdflatex guion_presentacion.tex  (dos pasadas)
% =============================================================================
\documentclass[11pt,a4paper]{article}

\usepackage[utf8]{inputenc}
\usepackage[T1]{fontenc}
\usepackage[spanish,es-nodecimaldot]{babel}
\usepackage{amsmath,amssymb}
\usepackage{booktabs}
\usepackage[margin=2.4cm]{geometry}
\usepackage[colorlinks=true,linkcolor=blue,urlcolor=blue]{hyperref}
\usepackage{enumitem}
\usepackage{xcolor}
\usepackage{mdframed}

\newcommand{\eps}{\varepsilon}
\newcommand{\gstar}{g^*}
\newcommand{\lammps}{\textsc{Lammps}}

% caja para las respuestas preparadas
\newmdenv[backgroundcolor=blue!5, linecolor=blue!40, linewidth=1pt,
          skipabove=4pt, skipbelow=8pt, innertopmargin=6pt,
          innerbottommargin=6pt]{respuesta}

\title{\textbf{Guion de estudio --- presentación a Gonzalo}\\
\large Partición de energía en el rebote de nanocúmulos LJ}
\author{Felipe Wasaff}
\date{Julio 2026}

\setlength{\parskip}{4pt}

\begin{document}
\maketitle

\begin{mdframed}[backgroundcolor=red!8, linecolor=red!60, linewidth=1.2pt]
\textbf{ADVERTENCIA (2026-07-19):} este guion corresponde al deck
\emph{antiguo} de 12 láminas, anterior al arbitraje simulado. La Parte~B
y partes de la Parte~A usan la física ya corregida
($N_c=N$, $\alpha_{\rm bola}\to1/2$, ``sin barras de error''), que la
versión v2 reemplazó. \textbf{Para la reunión con Gonzalo, estudiar
\texttt{borrador\_estudio\_reunion.pdf}} (sincronizado con las 16 láminas
de \texttt{presentacion\_gonzalo.pdf}). Este documento se conserva solo
por las explicaciones desde cero de la Parte~A y las respuestas C1--C12
que siguen vigentes.
\end{mdframed}

\noindent Este documento tiene tres partes: \textbf{(A)} los conceptos que
tienes que dominar antes de la reunión, explicados desde cero;
\textbf{(B)} el guion lámina por lámina (qué decir, qué señalar en cada
figura, cuánto tiempo); \textbf{(C)} las preguntas que Gonzalo
probablemente hará ---es profesor de mecánica estadística--- con las
respuestas preparadas.

\tableofcontents

% =============================================================================
\section{PARTE A --- Conceptos que debes dominar}
% =============================================================================

\subsection{Unidades de Lennard-Jones (por si pregunta por números)}

Todo el trabajo está en unidades reducidas: la longitud se mide en
$\sigma$ (el ``tamaño'' de un átomo, $\sim$3.4~\AA{} para argón), la
energía en $\eps$ (la profundidad del pozo atractivo entre dos átomos,
$\sim$0.0104~eV) y el tiempo en $\tau=\sqrt{m\sigma^2/\eps}\approx2.15$~ps.
Cuando digas ``la bola cae desde $26\,\sigma$'' estás diciendo ``desde
$\sim$9 nanómetros''; cuando digas ``el impacto dura $\sim2\,\tau$'' es
``$\sim$4 picosegundos''. La gravedad $\gstar=0.026\,\sigma/\tau^2$ es
\emph{enormemente} más grande que la gravedad real a esa escala --- es una
gravedad artificial para que el experimento quepa en tiempos de
simulación. Si Gonzalo pregunta: no afecta la física del impacto, porque
lo que importa es la velocidad de impacto $v_{\rm imp}=\sqrt{2gh}$, y esa
la controlamos con $h_0$.

\subsection{El coeficiente de restitución (COR) y por qué \texorpdfstring{$e<1$}{e<1}}

\begin{itemize}[leftmargin=1.4em]
    \item Definición: $e = |v_{\rm después}|/|v_{\rm antes}|$ del impacto.
    Como la altura de vuelo es $h=v^2/2g$, se cumple
    $\boxed{h_1/h_0=e^2}$.
    \item En los cursos de mecánica, $e$ se trata como constante del
    material. El primer resultado del trabajo es que \textbf{$e$ no es
    constante}: va de 0.92 (impacto suave) a 0.55--0.59 (impacto fuerte).
    Una pelota atómica rebota ``mejor'' cuando cae despacio.
    \item Intuición: a baja velocidad la bola no alcanza a excitar sus
    vibraciones internas (el impacto es ``adiabático'' respecto de los
    modos internos); a alta velocidad el golpe es violento comparado con
    el período de vibración de los átomos y la energía se reparte.
\end{itemize}

\subsection{Los cuatro canales --- la ecuación más importante del paper}

En el \textbf{primer apogeo} (el punto más alto después del primer
rebote, donde la velocidad del centro de masa es cero) se mide:
\begin{equation*}
    \alpha_{\rm cm} + \alpha_{T_B} + \alpha_{\Delta V_B}
    + \alpha_{\rm surf} = 1 \qquad \text{(todo dividido por $E_0$)}
\end{equation*}

\begin{enumerate}[leftmargin=1.6em]
    \item $\alpha_{\rm cm}=h_1/h_0^{\rm ef}=e_1^2$ --- la fracción que el
    centro de masa (CM) recuperó. Es la única energía ``útil'' para
    seguir rebotando.
    \item $\alpha_{T_B}$ --- energía cinética de los átomos de la bola
    \textsc{medida en el sistema del CM} (es decir, restando el
    movimiento de traslación). Es el ``calor'' cinético: los átomos
    vibran alrededor de sus posiciones. $T_B$ es literalmente lo que
    usarías para definir una temperatura interna instantánea de la bola
    ($T_B=(3N-3)\,k_B T_{\rm interna}/2$ si estuviera equilibrada).
    \item $\alpha_{\Delta V_B}$ --- el cambio de energía \textsc{potencial}
    interna de la bola entre antes del impacto y el apogeo. Ver
    \S\ref{sec:anarm}: es la ``expansión anarmónica''.
    \item $\alpha_{\rm surf}$ --- cuánta energía (cinética + potencial)
    ganó la superficie. Son los fonones: ondas elásticas que el golpe
    radia por el cristal.
\end{enumerate}

\textbf{El mensaje que tienes que poder decir en una frase:} ``la energía
que la pelota no recupera queda mayoritariamente \emph{dentro} de la
pelota ($\alpha_{T_B}+\alpha_{\Delta V_B}\approx0.4$--$0.6$), y el piso se
lleva poco (0.07--0.16). La intuición macroscópica de que `la energía se
disipa en el suelo' falla a esta escala.''

\subsection{La expansión anarmónica (\texorpdfstring{$\alpha_{\Delta V_B}$}{α\_ΔVB}) --- el canal ``nuevo''}
\label{sec:anarm}

Este es el canal conceptualmente más fino, prepáralo bien:

\begin{itemize}[leftmargin=1.4em]
    \item El potencial LJ \textbf{no es simétrico} alrededor del mínimo:
    es muy duro al comprimir ($r^{-12}$) y blando al estirar ($r^{-6}$).
    Eso es ``anarmonicidad''.
    \item Cuando la bola se calienta internamente ($T_B$ sube), sus
    átomos oscilan en ese pozo asimétrico y pasan \textbf{más tiempo del
    lado estirado}: la distancia media entre átomos crece. Es
    \textsc{exactamente} el mismo mecanismo de la \textbf{dilatación
    térmica} de los sólidos --- a Gonzalo le va a gustar esa conexión.
    \item Consecuencia: en el apogeo, la energía potencial interna $V_B$
    es \textsc{mayor} que antes del impacto (la bola está ``expandida''),
    y esa diferencia $\Delta V_B$ viaja con la bola. Es energía potencial
    almacenada, tan real como la cinética.
    \item Números: $\alpha_{\Delta V_B}\approx0.21$--$0.33$, comparable a
    $\alpha_{T_B}\approx0.22$. De hecho, que sean comparables \textsc{es}
    equipartición: en un oscilador, $\langle KE\rangle\approx\langle
    PE\rangle$ (virial). Si te preguntan por qué son parecidos, esa es
    la respuesta: ``es el teorema del virial actuando dentro de la
    bola''.
\end{itemize}

\subsection{¿Qué es el ``plateau''? (tu duda --- apréndete esto)}
\label{sec:plateau}

``Plateau'' = \emph{meseta}. Se refiere a esto: si graficas
$\alpha_{T_B}$ (la fracción cinética interna) contra la velocidad de
impacto $v_{\rm imp}$, la curva \textbf{sube y luego se aplana} en un
valor casi constante:

\begin{itemize}[leftmargin=1.4em]
    \item Para $v_{\rm imp}<0.9\,\sigma/\tau$: $\alpha_{T_B}$ es pequeño
    y crece con $v$ (régimen cuasi-elástico: el impacto apenas excita los
    modos internos).
    \item Para $v_{\rm imp}>0.9$--$1.1\,\sigma/\tau$: $\alpha_{T_B}$ se
    queda ``pegado'' en $\approx0.22$ sin importar cuánto más fuerte
    golpees. \textsc{Eso} es el plateau: la meseta de la curva
    $\alpha_{T_B}(v)$.
\end{itemize}

\textbf{Por qué existe la meseta (la física):} cuando el golpe es
suficientemente violento, el tiempo de contacto ($\sim2\,\tau$) basta
para que la energía se reparta (``termalice'') entre todos los grados de
libertad dinámicamente accesibles. Una vez que el reparto es completo, la
\textsc{fracción} que le toca a cada canal ya no depende de cuánta
energía entra --- solo de \textsc{cuántos} grados de libertad compiten.
Por eso la fracción se vuelve constante. Es el mismo motivo por el que el
calor específico de un sólido clásico no depende de la temperatura (ley
de Dulong-Petit): la equipartición reparte por igual entre modos, y la
fracción de cada modo es fija.

En cambio, si el golpe es suave ($v$ pequeña), la termalización \textsc{no}
alcanza a ocurrir y la fracción sí depende de $v$: por eso la curva sube
antes de aplanarse. Y en el cúmulo de $N=13$, que es muy rígido (una sola
capa cerrada, frecuencias propias altas), la termalización no ocurre ni
siquiera a $v=2$: \textbf{$N=13$ nunca alcanza el plateau} en el rango
estudiado. Por eso en la figura va con símbolo abierto y excluido del
ajuste.

\subsection{La fórmula de equipartición y \texorpdfstring{$N_c$}{Nc}}
\label{sec:equip}

\begin{equation*}
    \alpha_{\rm int}(N) = \frac{3N-3}{3N-3+3N_c}
\end{equation*}

\begin{itemize}[leftmargin=1.4em]
    \item \textbf{De dónde sale:} si en el contacto la energía $E_0$ se
    reparte por igual entre $f_{\rm bola}=3N-3$ grados de libertad
    internos de la bola y $f_{\rm sup}=3N_c$ grados de libertad
    ``activos'' de la superficie, la fracción de la bola es
    $f_{\rm bola}/(f_{\rm bola}+f_{\rm sup})$. Eso es todo: conteo de
    grados de libertad, puro espíritu de equipartición.
    \item \textbf{Por qué $3N-3$ y no $3N-6$:} un cúmulo de $N$ átomos
    tiene $3N$ grados de libertad; se restan 3 de traslación del CM
    (esos son $\alpha_{\rm cm}$, no son ``internos''). Las 3 rotaciones
    \textsc{sí} se cuentan como internas porque un impacto no
    perfectamente centrado excita rotación. Para $N$ grande la diferencia
    entre $3N-3$ y $3N-6$ es irrelevante; para $N=2$ o 7 no lo es, y por
    eso esos tamaños son solo controles cualitativos.
    \item \textbf{$N_c$ = número efectivo de átomos de la superficie que
    participan.} No es un parámetro microscópico medible directamente; se
    \textsc{calibra} con los datos. El resultado interesante: si usas la
    fracción interna \textsc{total}
    ($\alpha_{\rm bola}=\alpha_{T_B}+\alpha_{\Delta V_B}$), la
    calibración da $\boxed{N_c\approx N}$. Lectura física: cada átomo de
    la bola ``recluta'' aproximadamente un átomo del piso. Plausible
    porque la bola de radio $\sim N^{1/3}$ penetra una profundidad
    comparable a su radio: volumen activado $\sim N$, no solo el área de
    contacto $\sim N^{2/3}$.
    \item \textbf{La predicción estrella:} con $N_c=N$, la fórmula da
    $\alpha_{\rm bola}\to 3N/(3N+3N)=\boxed{1/2}$ cuando $N\to\infty$.
    Una nanopelota grande, fría y sin adhesión absorbe como energía
    interna la \textsc{mitad} de su energía de impacto. Límite limpio,
    memorizable, falsable.
    \item \textbf{La salvedad honesta (dila tú antes de que la encuentre
    Gonzalo):} los datos de $N=79$ y 87 dan $\alpha_{\rm bola}=0.543$,
    \textsc{levemente por encima} de la cota $1/2$. Estrictamente eso
    implica $N_c<N$ (invertir la fórmula da $N_c/N=1.09$, $0.83$, $0.83$
    para $N=55$, 79, 87). Origen probable: (i)~la energía rotacional está
    incluida en $\alpha_{\rm bola}$ pero la rotación no es un oscilador
    (no tiene PE asociada, ``cuenta distinto'' en equipartición), o
    (ii)~$\alpha_{\rm surf}$ está algo subestimado. Distinguirlas
    requiere medir el momento angular post-rebote (\texttt{compute
    angmom} en \lammps) --- está declarado como trabajo pendiente.
\end{itemize}

\subsection{\texorpdfstring{$h^*=T_B/(Mg)$}{h*=TB/(Mg)}: qué es y qué NO es}

\begin{itemize}[leftmargin=1.4em]
    \item Idea: si al primer rebote $\alpha_{\rm cm}\approx\alpha_{T_B}$
    (mira la tabla: 0.304 vs 0.319 en la fila $h_0^{\rm ef}=25.9$),
    entonces la altura del primer rebote
    $h_1=\alpha_{\rm cm}h_0^{\rm ef}\approx\alpha_{T_B}h_0^{\rm ef}
    =T_B/(Mg)$. Es una \textsc{coincidencia numérica} del primer rebote
    elevada a escala útil: conociendo el calentamiento interno predices
    la altura. Funciona: $h_1/h^*=0.96$.
    \item Lo que \textsc{no} es: un \textbf{atractor}. ``Atractor'' =
    estado al que la dinámica converge y en el que se queda. Uno podría
    conjeturar que las alturas de rebote convergen a $h^*(T_B)$ porque
    $T_B$ ``recarga'' el rebote. La simulación larga ($5000\,\tau$, 23
    rebotes) muestra lo contrario: la razón $h_n/h_n^*$ cae de 0.96 a
    $\sim$0.18. La altura decae geométricamente (cada rebote pierde un
    factor $e^2$) mientras que $T_B$, una vez depositada, casi no cambia
    (baja solo $\times0.7$ en 23 rebotes). Son dos relojes distintos:
    $h_n$ decae rápido, $T_B$ lento. Al final la bola queda
    \textsc{atrapada} vibrando a menos de $1\,\sigma$ del piso --- el
    análogo exacto del ``prrrrt'' final de una pelota real.
    \item El ``colapso'' en $t\approx713\,\tau$ (la bola se cuela
    \textsc{entre} los átomos del piso) es un \textbf{artefacto}: con
    interacción puramente repulsiva y sin atracción lateral, los átomos
    de la bola terminan filtrándose por los huecos de la red. Lo decimos
    explícitamente en el paper. Saber distinguir física de artefacto es
    parte del mensaje pedagógico.
\end{itemize}

\subsection{Adhesión (\texorpdfstring{$c_{ps}=1$}{cps=1})}

\begin{itemize}[leftmargin=1.4em]
    \item Restauramos el término atractivo bola--piso del LJ. Resultado:
    en los 24 casos (4 tamaños $\times$ 6 alturas) la bola queda pegada
    al primer contacto. No es sorpresa: la energía de adhesión medida
    ($E_{\rm adh}=52$--$209\,\eps$ según $N$) es mucho mayor que la
    energía CM disponible tras el rebote
    ($\alpha_{\rm cm}E_0\sim5$--$50\,\eps$).
    \item $E_{\rm adh}\propto N^{0.73}\approx N^{2/3}$ = escalado del
    \textsc{área} de contacto (los átomos del interior no tocan el piso).
    Compara con la masa $\propto N$. De ahí la predicción
    contraintuitiva: el umbral de caída para lograr rebote escala
    $h_0^*\propto N^{0.73}/N=N^{-0.27}$ --- \textbf{los cúmulos grandes
    escapan más fácil} (más energía por unidad de adhesión).
    \item El valor absoluto del umbral ($\sim10^2$--$10^3\,\sigma$) es
    \textsc{solo} orden de magnitud: a esas velocidades la bola se funde
    al impactar (lo vio el propio trabajo 2012 vía $g(r)$), el área de
    contacto crecería con $v$, y $\alpha_{\rm cm}$ está tomado del caso
    sin adhesión. Preséntalo siempre con esa triple salvedad.
    \item Conexión con literatura: \emph{soft landing} de cúmulos sobre
    superficies (deposición de nanopartículas sin destruirlas) ---
    Awasthi et al., Phys.\ Rev.\ B.
\end{itemize}

\subsection{Los detalles técnicos que dan credibilidad}
\label{sec:tecnicos}

\begin{enumerate}[leftmargin=1.6em]
    \item \textbf{Cúmulos de $N$ exacto}: en 2012 la bola se creaba
    recortando una esfera de una red $\to$ $N$ variable (79--84) y bola
    no relajada. Ahora: capas FCC completas (13, 55, 79, 87\ldots)
    minimizadas en vacío $\to$ $N$ exacto, bola inicialmente fría. Sin
    esto, la medición de $\alpha_{T_B}$ estaría contaminada con energía
    interna inicial espuria.
    \item \textbf{Potencial truncado y desplazado (shift)}: el LJ se
    corta en $r_c=2.5\,\sigma$. Si solo lo cortas,
    $U(r_c)=-0.0163\,\eps\neq0$, y cada par de átomos que cruza $r_c$
    hace saltar la energía potencial \textsc{registrada} en
    $\pm0.0163\,\eps$. Con miles de cruces por impacto, el balance de
    energía fallaba en 2--6\% de $E_0$ --- lo descubrimos porque el
    cierre $\sum\alpha_i=1$ no cuadraba al medir los 4 canales de forma
    independiente. La solución estándar: desplazar el potencial
    ($U-U(r_c)$), que \textsc{no} cambia las fuerzas (la derivada es la
    misma) y por tanto \textsc{no} cambia ninguna trayectoria --- solo
    arregla la contabilidad. Con shift: conservación a
    $|\Delta E|<0.03\,\eps\approx4\times10^{-4}E_0$ por impacto y cierre
    $<10^{-4}$. Detalle importante: los resultados anteriores eran
    válidos porque $\alpha_{\Delta V_B}$ se calculaba por residuo.
    \item \textbf{Digitalización de la Fig.~7 de 2012}: extrajimos las
    curvas del PDF del capítulo por color de píxel, calibrando con los 6
    valores conocidos del rebote 0 (error rms de calibración:
    $0.08\,\sigma$). Hallazgo: en 2012 $h_0$ es la \textsc{altura
    absoluta} del CM (las curvas convergen al nivel de reposo
    $\approx10.3\,\sigma$), mientras que nuestro $h_0$ se mide desde el
    tope de la caja. Con la \textsc{misma} etiqueta $h_0$, las caídas
    efectivas difieren en $\sim8\,\sigma$. Comparando a caída efectiva
    igual: acuerdo excelente a baja energía ($e_1$: 0.90 de ellos vs
    0.92 nuestro) y 2012 queda $\sim$0.06--0.09 abajo a alta energía
    (explicable: bola sin minimizar + $N$ variable).
    \item \textbf{Reproducibilidad}: cada figura del paper la genera un
    script del repo desde los datos crudos. Cualquier número que Gonzalo
    cuestione, se puede regenerar delante de él.
\end{enumerate}

% =============================================================================
\section{PARTE B --- Guion lámina por lámina ($\sim$15 min total)}
% =============================================================================

\subsection*{Lámina 1 --- Portada (15 s)}
``Te traigo el estado del trabajo del rebote: qué reprodujimos del
capítulo 2012, qué encontramos nuevo, y una propuesta concreta de
publicación que quiero cerrar contigo hoy.''

\subsection*{Lámina 2 --- De dónde venimos (1 min)}
\begin{itemize}[leftmargin=1.4em]
    \item Señala el bloque de la derecha: ``la pregunta que ordena todo
    el paper es esta: $e<1$ te dice \textsc{cuánta} energía no vuelve;
    nosotros respondemos \textsc{dónde} queda.''
    \item Recorre los 5 bullets de la izquierda rápido; no te detengas en
    detalles técnicos (vuelven en la lámina 10).
    \item Frase puente: ``el resultado central cabe en una imagen --- es
    la siguiente.''
\end{itemize}

\subsection*{Lámina 3 --- La imagen OVITO (2 min; deja que la mire)}
\begin{itemize}[leftmargin=1.4em]
    \item Deja 10 segundos de silencio; la figura se explica sola.
    \item Recorrido: ``(a) cae --- toda la energía es mecánica del CM.
    (b) compresión máxima --- mira la barra: la energía está como
    deformación, 54\% potencial intra-bola, y el canal gris $U_c$ es el
    contacto, que solo existe durante la colisión. (c) al separarse ya
    está caliente --- mira el moteado de colores en la bola. (d) apogeo:
    31\% cinética interna, 31\% potencial interna, 6\% superficie, 32\%
    altura recuperada.''
    \item Remate: ``o sea: la bola se lleva \textsc{adentro} el 62\% de
    la energía; el piso se queda con el 6\%. La intuición de `se disipa
    en el suelo' falla.''
    \item Detalle si pregunta por los colores: bola coloreada por KE
    interna (sin el CM), superficie por KE con escala 30 veces más
    sensible para ver los fonones.
\end{itemize}

\subsection*{Lámina 4 --- Los cuatro canales (2 min)}
\begin{itemize}[leftmargin=1.4em]
    \item ``Esto es lo mismo cuantificado para seis alturas de caída.''
    \item Señala la columna $\alpha_{T_B}$: ``aquí está el plateau: sobre
    $v\approx0.9$, la fracción cinética interna se clava en 0.21--0.24.
    La fracción deja de depender de cuán fuerte golpeas --- señal de que
    el contacto alcanza a termalizar'' (si pide detalle $\to$
    \S\ref{sec:plateau}).
    \item Señala $\alpha_{\Delta V_B}$: ``y este es el canal que nadie
    había separado: la expansión anarmónica --- dilatación térmica
    instantánea de la bola. Es tan grande como el cinético; es el virial
    actuando.''
    \item Señala la fila 1: ``impacto lento = casi elástico, 85\%
    recuperado.''
    \item Cierre con el bloque rojo: ``suma exacta 1, cierre medido mejor
    que $10^{-4}$ --- la contabilidad es completa.''
\end{itemize}

\subsection*{Lámina 5 --- Validación contra 2012 (2 min; lámina delicada, tono de colaboración)}
\begin{itemize}[leftmargin=1.4em]
    \item ``Antes de extender, validamos contra ustedes. Como el capítulo
    no publica $e_1$ numérico, digitalicé la Fig.~7 por color, calibrando
    con las alturas iniciales conocidas.''
    \item El hallazgo de convención: ``encontré que en el capítulo $h_0$
    es la altura absoluta del CM --- las curvas parten en $h_0$ y
    descansan en $\sim10.3\,\sigma$. La nuestra se mide desde el tope de
    la caja. Misma etiqueta, caída distinta; hay que comparar a caída
    efectiva igual.''
    \item ``Así comparado: 0.90 ustedes, 0.92 nosotros en el régimen
    suave --- excelente. A alta energía ustedes quedan un poco abajo,
    consistente con la bola sin minimizar y el $N$ variable del
    \texttt{create\_atoms}.''
    \item \textsc{No} digas ``error de ustedes''; di ``diferencia de
    convención documentada''.
\end{itemize}

\subsection*{Lámina 6 --- Atribución (2 min; LA lámina de la conversación)}
\begin{itemize}[leftmargin=1.4em]
    \item Léela casi literal: ``esto lo necesito confirmar contigo antes
    de someter. El capítulo establece los canales, las alturas por
    rebote, el rol de $U_c$ durante el contacto y la fusión. El borrador
    presenta como contribución nueva: el plateau cuantificado
    $\alpha_{T_B}\approx0.22$, la fórmula de equipartición con $N_c$, y
    la escala $h^*$.''
    \item Pregunta directa: ``¿esto coincide con tu recuerdo? ¿Algo de
    esto viene de notas tuyas de la época que debamos citar?''
    \item Escucha. Si algo era de él, la solución es simple y ya está
    prevista: se cita como comunicación privada o pasa a la introducción
    como idea previa. No es un conflicto, es limpieza de autoría.
\end{itemize}

\subsection*{Lámina 7 --- Equipartición (2.5 min; la lámina más ``mecánica estadística'')}
\begin{itemize}[leftmargin=1.4em]
    \item Panel (a): ``cada curva es un tamaño de bola. Los grandes (55,
    79, 87) suben y se aplanan --- ese es el plateau. El de 13 (cuadrados
    abiertos azules) nunca se aplana: demasiado rígido para termalizar en
    $2\,\tau$ de contacto. Se comporta como molécula, no como sólido.''
    \item Panel (b): ``el valor del plateau contra $N$. La fórmula es
    puro conteo de grados de libertad: $3N-3$ de la bola contra $3N_c$
    del piso.''
    \item El límite: ``calibrando con $\alpha_{\rm bola}$ total sale
    $N_c\approx N$ --- cada átomo de la bola recluta uno del piso --- y
    eso da el límite $\alpha_{\rm bola}\to1/2$.''
    \item La salvedad \textsc{antes} de que la encuentre él: ``los datos
    de 79 y 87 dan 0.543, un pelo sobre $1/2$. Lo decimos en el paper:
    probablemente la rotación, que está en $\alpha_{\rm bola}$ pero no es
    un oscilador. Medir el momento angular es el siguiente paso.''
\end{itemize}

\subsection*{Lámina 8 --- \texorpdfstring{$h^*$}{h*} (1.5 min)}
\begin{itemize}[leftmargin=1.4em]
    \item ``Del plateau sale una escala: $h^*=T_B/Mg$. Primer rebote:
    0.96 --- funciona. ¿Es un atractor? No: la razón cae a 0.18 porque la
    altura decae geométricamente y $T_B$ casi no cambia. La bola termina
    atrapada a menos de $1\,\sigma$ tras 23 rebotes.''
    \item Artefacto: ``el desplome final a través de la red es un
    artefacto del potencial repulsivo puro y está documentado como tal.''
\end{itemize}

\subsection*{Lámina 9 --- Adhesión (1.5 min)}
\begin{itemize}[leftmargin=1.4em]
    \item ``Con la atracción encendida, pega siempre: 24 de 24.''
    \item ``$E_{\rm adh}$ escala como $N^{0.73}\approx$ área de contacto;
    la masa como $N$. Por eso los grandes escapan más fácil: umbral
    $\propto N^{-0.27}$.''
    \item ``El valor absoluto del umbral es solo orden de magnitud --- a
    esas velocidades la bola se funde, como mostraron ustedes con
    $g(r)$.''
\end{itemize}

\subsection*{Lámina 10 --- Rigor y límites (1.5 min)}
\begin{itemize}[leftmargin=1.4em]
    \item Columna izquierda, el shift en una frase: ``descubrimos que sin
    desplazar el potencial truncado, la contabilidad energética falla en
    2--6\% por impacto --- puro artefacto de \emph{bookkeeping}, las
    trayectorias no cambian. Con shift, cierre exacto.''
    \item Columna derecha, sé el primero en decirlo: ``$T=0$, una
    orientación, una trayectoria por punto: aún no hay barras de error
    estadísticas. Eso define la extensión natural del trabajo, y está
    declarado en el paper.''
\end{itemize}

\subsection*{Lámina 11 --- Plan editorial (1.5 min)}
\begin{itemize}[leftmargin=1.4em]
    \item ``Según lo que conversaron, apunté a la RBEF, sección Artigos
    Gerais: física computacional con punto de vista didáctico. El
    manuscrito ya cumple el protocolo: bilingüe, referencias en orden,
    formato de la revista, y agregué una sección de actividades para el
    aula.''
    \item ``Además existe la versión de investigación en doble columna
    (español e inglés) si después queremos un segundo paper con la
    estadística completa.''
    \item Flujo: tu revisión $\to$ arXiv (\texttt{physics.ed-ph}) $\to$
    submissão $\to$ repo público.
\end{itemize}

\subsection*{Lámina 12 --- Decisiones (2 min; cierre operativo)}
Recorre las 5 y anota sus respuestas:
\begin{enumerate}[leftmargin=1.6em]
    \item Atribución (viene de la lámina 6).
    \item Idioma del envío (español está aceptado por la RBEF).
    \item Su ORCID + visto bueno a arXiv.
    \item Autor correspondiente y orden de autores.
    \item ¿Guardamos la versión PRE para más adelante con la estadística?
\end{enumerate}

% =============================================================================
\section{PARTE C --- Preguntas probables de Gonzalo (y tus respuestas)}
% =============================================================================

\subsection*{C1. ``¿Equipartición? Esto no está en equilibrio. ¿Con qué
derecho usan el teorema?'' (la más probable de un mecánico estadístico)}
\begin{respuesta}
Correcto, no hay equilibrio ni ensamble --- hay \textsc{una} colisión.
Usamos equipartición como \emph{ansatz de reparto}: si el tiempo de
contacto basta para redistribuir la energía entre los modos dinámicamente
accesibles, cada uno recibe en promedio lo mismo. No lo suponemos: lo
\textsc{contrastamos}. La evidencia de que el reparto ocurre es el
plateau (la fracción deja de depender de la energía inyectada) y de que
\textsc{no} ocurre cuando no puede ocurrir: $N=13$, rígido, no plateau.
La fórmula es fenomenológica con un parámetro calibrado ($N_c$); su
mérito es predecir el \textsc{escalado} con $N$ y el límite $1/2$.
\end{respuesta}

\subsection*{C2. ``¿Y por qué la energía cinética interna equivale a la
mitad de la interna total? ¿Virial?''}
\begin{respuesta}
Sí: para osciladores, $\langle KE\rangle=\langle PE\rangle$. Por eso
$\alpha_{T_B}\approx\alpha_{\Delta V_B}$ en el plateau (0.22 vs
0.21--0.33; la PE fluctúa más porque incluye la deformación anarmónica
neta, no solo la oscilación).
\end{respuesta}

\subsection*{C3. ``¿Cómo definen temperatura de la bola si son 79
átomos?''}
\begin{respuesta}
No la definimos: reportamos $T_B$, la KE interna, que es una cantidad
mecánica bien definida sin hipótesis de equilibrio. Si uno \textsc{quiere}
una temperatura instantánea, $T=2T_B/(3N-3)k_B$, pero el paper trabaja
siempre con fracciones de energía, nunca con temperatura.
\end{respuesta}

\subsection*{C4. ``$3N-3$ y no $3N-6$, ¿por qué cuentan la rotación como
modo interno?''}
\begin{respuesta}
Porque un impacto no perfectamente central transfiere momento angular;
esa energía no está en el CM ni en la superficie, queda ``en la bola''.
La contamos como interna, y decimos explícitamente que la rotación no es
un oscilador --- probablemente explica el pequeño exceso sobre $1/2$.
Para $N\geq13$ el cambio $3N-3\to3N-6$ mueve la fórmula menos que las
otras incertezas.
\end{respuesta}

\subsection*{C5. ``El $e_1$ no es monótono con $h_0$ (0.55 en
$25.9\,\sigma$ pero 0.63 en $41.8\,\sigma$). ¿Eso es física o ruido?''}
\begin{respuesta}
No lo sabemos aún, y lo decimos: con $T=0$ y una sola orientación hay
\textsc{una} trayectoria por punto, sin estadística. Puede ser resonancia
con la fase de la vibración interna al llegar al piso, o varianza de
condición inicial. Es la primera prioridad del trabajo futuro (3--5
réplicas con orientaciones aleatorias). Nota: la serie de ustedes en 2012
también tiene un no-monótono en esa zona ($e_1=0.43$ en su tercera
altura).
\end{respuesta}

\subsection*{C6. ``¿La superficie no es demasiado chica? ¿El fonón no
rebota en el borde y vuelve?''}
\begin{respuesta}
La losa es $8\times8\times6$ celdas con las 2 capas de abajo congeladas y
bordes periódicos laterales. En el tiempo hasta el primer apogeo
($\sim25\,\tau$) la onda alcanza a recorrer la caja, así que
$\alpha_{\rm surf}$ incluye energía que re-circula por los bordes
periódicos --- posible subestimación del flujo saliente ``verdadero'' de
una superficie semi-infinita. Es una limitación conocida de las MD de
impacto; agrandar la losa y medir $\alpha_{\rm surf}(L)$ es un chequeo
pendiente razonable. [Sé honesto: no lo hemos barrido.]
\end{respuesta}

\subsection*{C7. ``¿Termostato? ¿Cómo controlan $T$?''}
\begin{respuesta}
No hay termostato: NVE puro, $T$ inicial $=0$ exacto (velocidades cero,
cúmulo minimizado). Toda la energía viene de la caída. Por eso la
conservación de energía es \textsc{el} control de calidad, y por eso el
asunto del shift importaba tanto.
\end{respuesta}

\subsection*{C8. ``¿Paso de tiempo? ¿Se conserva la energía en el
impacto?''}
\begin{respuesta}
$\Delta t=5\times10^{-3}\,\tau$, velocity-Verlet. Validado midiendo la
energía total (incluida la PE gravitatoria) a lo largo de un impacto
completo: $|\Delta E|<0.03\,\eps\approx4\times10^{-4}E_0$. Ustedes usaron
Beeman con $\Delta t$ 10 veces menor; con Verlet y potencial desplazado
no hace falta.
\end{respuesta}

\subsection*{C9. ``¿Por qué la superficie se lleva tan poco? Uno
esperaría acoplo fuerte al continuo de fonones.''}
\begin{respuesta}
Dos razones: (i)~impedancia --- el contacto dura $\sim2\,\tau$ y excita
preferentemente los modos de la huella, no el continuo; (ii)~el canal
dominante durante el contacto ($U_c$) es potencial de contacto, y esa
energía \textsc{regresa} mayoritariamente a la bola al separarse (es lo
que se ve en la barra gris de la lámina 3). El análisis de 2012, hecho
durante el contacto, veía $U_c$ grande; nuestra partición al apogeo
muestra su destino final. Ambas descripciones son compatibles: miden en
instantes distintos.
\end{respuesta}

\subsection*{C10. ``¿Y esto es publicable en RBEF? Es una revista de
enseñanza.''}
\begin{respuesta}
Sección ``Artigos Gerais'': acepta física
teórica/computacional/experimental ``bajo un punto de vista didáctico,
con temas poco discutidos en textos clásicos, a nivel accesible a no
especialistas'' --- es exactamente este manuscrito. Sin costo de
publicación, SciELO, acepta español. Y la versión de investigación queda
disponible para una revista de física computacional cuando tengamos la
estadística de orientaciones (ahí sí PRE o EPJ~B).
\end{respuesta}

\subsection*{C11. ``¿Qué pasa con $N$ intermedios? ¿19, 43?''}
\begin{respuesta}
Son las capas completas que faltan entre 13 y 55, justo donde ocurre la
transición molécula$\to$sólido. Están en la lista de trabajo futuro; los
cúmulos son triviales de generar con el mismo script.
\end{respuesta}

\subsection*{C12. ``¿El $N=2$ con $\alpha_{\rm surf}=0.18$? ¿Dónde se fue
el resto?''}
\begin{respuesta}
$e_1=0.90\to\alpha_{\rm cm}=0.81$; $\alpha_{\rm surf}=0.18$;
$\alpha_{\rm bola}\approx0.003$. Suma $\approx1$. El dímero vertical casi
no acopla el golpe a su único modo de vibración (que es perpendicular al
movimiento si cae ``acostado''): todo lo que no vuelve al CM se lo lleva
el piso. Es el único caso donde la intuición macroscópica (``el piso
absorbe'') es correcta --- porque la bola no tiene dónde guardar energía.
\end{respuesta}

% =============================================================================
\section*{Chuleta de números (para tener a mano)}
% =============================================================================

\begin{center}
\begin{tabular}{ll}
\toprule
\textbf{Cantidad} & \textbf{Valor} \\
\midrule
Plateau $\alpha_{T_B}$ ($N=79$) & 0.21--0.24, arranca en $v\approx0.9\,\sigma/\tau$ \\
$\alpha_{\Delta V_B}$ & 0.21--0.33 \\
$\alpha_{\rm surf}$ & 0.07--0.16 \\
$\alpha_{\rm bola}$ ($N=79,87$, plateau) & 0.543 (cota teórica $1/2$) \\
$N_c$ calibrado & $\approx N$ (1.09 / 0.83 / 0.83 para 55/79/87) \\
$h_1/h^*$ primer rebote & 0.96; cae a $\sim$0.18 en $n=20$--23 \\
Atrapamiento & 23 rebotes, $t\approx660\,\tau$; colapso (artefacto) $t\approx713\,\tau$ \\
$E_{\rm adh}$ & 52--209$\,\eps$; $\propto N^{0.73}$; umbral $\propto N^{-0.27}$ \\
Conservación con shift & $|\Delta E|<0.03\,\eps\approx4\times10^{-4}E_0$; cierre $<10^{-4}$ \\
Validación vs 2012 & $e_1$ 0.90 vs 0.92 (baja $E$); $-0.06$ a $-0.09$ (alta $E$) \\
Costo & $\sim$1.5 min/rebote en notebook, 1 núcleo \\
\bottomrule
\end{tabular}
\end{center}

\end{document}
